\chapter[Introdução]{Introdução}
\label{cap:introducao}
Este capítulo introduz o(a) leitor(a) aos principais temas deste trabalho. São contextualizados o uso de dados na administração pública federal e o papel de plataformas governamentais, como o Gov Hub, na consolidação desses dados, bem como a passagem de análises descritivas para análises preditivas apoiadas por aprendizado de máquina e a disciplina  de MLOps. Em seguida, são apresentados o problema da paralisação de obras do Programa Minha Casa, Minha Vida e a justificativa para integrar o ciclo de vida de modelos a uma plataforma de dados existente. Na sequência, são expostos a questão de pesquisa que guia toda esta investigação e os objetivos. Por fim, são descritas a metodologia em síntese e a organização dos demais capítulos.
\section{Contextualização}

A administração pública federal brasileira produz, em seus sistemas
estruturantes, um volume crescente de dados sobre a execução de políticas
públicas. A Política de Dados Abertos do Poder Executivo federal
\cite{brasil2016dadosabertos} e a Lei do Governo Digital
\cite{brasil2021governodigital} estabeleceram a disponibilização e o uso
desses dados como diretriz de gestão. Disponibilizar o dado, porém, não
garante que ele apoie decisões: \citeonline{janssen2012} mostram que os
benefícios dos dados abertos dependem de capacidades institucionais para
integrá-los, interpretá-los e mantê-los ao longo do tempo.

Plataformas de dados governamentais como o Gov Hub, desenvolvido pelo
Laboratório de Competência em Software Livre (LabLivre) da Universidade de
Brasília em parceria com órgãos federais, respondem a parte desse desafio. Elas ingerem dados de
sistemas heterogêneos, os organizam em camadas de qualidade crescente e os
publicam em painéis analíticos. No caso do Ministério das Cidades (MCID), a
plataforma consolida o histórico mensal de execução das obras do Programa
Minha Casa, Minha Vida (MCMV) \cite{brasil2023mcmv}, com situação, percentuais
de execução física e financeira, valores e datas de cada empreendimento.

Com esse histórico disponível, a pergunta deixa de ser descritiva (``quais
obras estão paralisadas?'') e passa a ser preditiva (``quais obras ativas vão
paralisar nos próximos meses?''). Responder a ela exige modelos de
aprendizado de máquina (\textit{Machine Learning}, ML). Colocar esses modelos
em produção de forma confiável é um problema de engenharia distinto, conhecido
como MLOps (\textit{Machine Learning Operations})
\cite{kreuzberger2023}.

\section{Problema e justificativa}

A paralisação de obras públicas é um problema de grande escala no Brasil.
Auditoria operacional do Tribunal de Contas da União identificou 14.403 obras
paralisadas em um universo de 38.412 obras financiadas com recursos federais
\cite{tcu2019}. Quando a fiscalização só identifica a paralisação depois que
ela ocorre, o custo de retomada já incorreu.

No Gov Hub Cidades, a identificação de obras em risco é feita hoje por uma
regra fixa: uma obra é sinalizada quando passa 90 dias sem atualização,
apresenta diferença superior a 10 pontos percentuais entre execução física e
financeira ou tem menos de 30\% de execução. Essa regra nunca foi comparada
com as paralisações que de fato ocorreram. Um modelo preditivo treinado no
histórico poderia, em princípio, antecipar mais paralisações com o mesmo
número de vistorias.

O obstáculo não é apenas estatístico. A literatura mostra que o código de
treinamento de um modelo é uma fração pequena de um sistema de ML em
produção, e que a maior parte do esforço e da dívida técnica está na
infraestrutura em torno dele: dependências de dados, pipelines, configuração e
monitoramento \cite{sculley2015}. Estudos de casos reais apontam desafios em
todas as etapas da implantação \cite{paleyes2022}, e no setor público eles se
somam a obstáculos institucionais, como acesso restrito a dados, conjuntos de
dados sem versionamento e governança frágil de proveniência
\cite{ferreira2025}. Em plataformas governamentais que já adotam práticas
maduras de engenharia de dados, como um \textit{Data Lakehouse} organizado
nas camadas bronze, prata e ouro da arquitetura medalhão, orquestração de
fluxos, transformação declarativa e catálogo de dados, falta definir como o
ciclo de vida de modelos se encaixa nessas práticas sem duplicar lógica nem
criar camadas paralelas.

\section{Questão de pesquisa}

Diante desse cenário, este trabalho é orientado pela seguinte questão de
pesquisa:

\begin{quote}
\textbf{QP:} Como estruturar uma arquitetura de MLOps que se integre a uma
plataforma de dados governamentais existente, garantindo reprodutibilidade,
rastreabilidade e avaliação honesta de modelos preditivos?
\end{quote}

Dela derivam três questões específicas:

\begin{itemize}
	\item \textbf{QE1:} Quais responsabilidades do ciclo de vida de ML devem
	ficar na camada de transformação de dados, quais no ambiente de modelagem e
	quais no orquestrador, e em qual camada do \textit{lakehouse} fica cada
	artefato de ML?
	\item \textbf{QE2:} Como garantir que qualquer resultado de um modelo possa
	ser reproduzido a partir do código e dos dados que o geraram?
	\item \textbf{QE3:} A arquitetura proposta permite avaliar um modelo contra
	a regra vigente de forma comparável, sem vazamento de informação temporal?
\end{itemize}

\section{Objetivos}

\subsection{Objetivo geral}

Propor, implementar e avaliar uma arquitetura de MLOps para plataformas de
dados governamentais, instanciada no caso de previsão de paralisação de obras
do Programa Minha Casa, Minha Vida na plataforma Gov Hub Cidades.

\subsection{Objetivos específicos}

\begin{enumerate}
	\item Levantar, na literatura, os princípios, componentes e riscos de
	sistemas de ML em produção, com atenção ao contexto do setor público;
	\item Definir os requisitos da arquitetura a partir das restrições da
	plataforma de dados existente e do caso de uso;
	\item Projetar a arquitetura sobre as camadas do \textit{lakehouse},
	especificando os papéis que a compõem e as responsabilidades de cada um,
	independentemente de ferramenta;
	\item Instanciar a arquitetura no caso de prevenção de paralisação de
	obras, cobrindo preparação de dados, treino, avaliação, escore mensal e
	monitoramento;
	\item Avaliar a arquitetura quanto à reprodutibilidade, à rastreabilidade e
	à capacidade de comparar o modelo com a regra vigente.
\end{enumerate}

\section{Metodologia em síntese}

O trabalho segue a \textit{Design Science Research} (DSR)
\cite{hevner2004,peffers2007}, abordagem voltada à construção e avaliação de
artefatos que resolvem problemas reais. O artefato principal é a arquitetura
de MLOps. O caso de previsão de paralisação de obras serve como demonstração
e avaliação em contexto real. O Capítulo~\ref{cap:metodologia} detalha as
etapas.

\section{Organização do trabalho}

O restante deste documento está organizado da seguinte forma. O
Capítulo~\ref{cap:referencial} apresenta o referencial teórico sobre
plataformas de dados, sistemas de ML em produção, MLOps e o uso de ML no
setor público, além dos trabalhos relacionados. O
Capítulo~\ref{cap:metodologia} descreve a metodologia de pesquisa. O
Capítulo~\ref{cap:proposta} apresenta a arquitetura proposta e sua
instanciação no caso das obras do MCMV. O Capítulo~\ref{cap:resultados}
relata os resultados preliminares obtidos até esta etapa. Por fim, o
Capítulo~\ref{cap:consideracoes} traz as considerações finais e o cronograma
para a continuidade do trabalho.
